\noindent\textit{2015 manuscript.}

\section{Heating and the dilute gas}
\noindent\textit{Manuscript section title: Background.}
\cite{RBaierlein1999}
\subsection{Heating and Temperature}

Heating: keep in mind 3 different types of heating for energy exchange between two systems:

\begin{enumerate}
  \item Heating by conduction - literal contact, molecules jiggle faster from molecules jiggling faster by bouncing off them
  \item Heating by radiation - em waves from hot source strike and excite target
  \item heating by convection - energy transport by flow (perhaps a fluid)
\end{enumerate}

This all relates to \\
$Q$

\subsection{Some dilute gas relationships}
\label{sec:dilute-gas}

\subsubsection*{Pressure according to kinetic theory} i.e. some kinetic theory

$F \equiv $ force on area $A$ due to molecules \\
$\Delta p \equiv $ momentum transferred to wall per collision \\
$n \equiv$ number of collisions in time $\Delta t$

Thus
\[
F = \frac{ (\Delta p)n }{ \Delta t }
\]

Now $\Delta p = 2 m v_x$ since $\Delta p = mv_x - (-mv_x) = 2mv_x$ (elastic collision with momentum conversation)

$v_x \Delta t A$ is a volume, inside of which gas molecules can be within distance $v_x \Delta t$ toward the wall.  \\
$\frac{N}{V}$ number density of molecules

Assume equal distribution of velocities: Thus $\frac{1}{2}$
\[
 n = (v_x \Delta t A) \frac{1}{2} \frac{N}{V}
\]
\[
P = \frac{F}{A} = \frac{ (2mv_x)(v_x \Delta t A ) \frac{1}{2} \frac{N}{V} }{ A \Delta t} = mv_x^2 \frac{N}{V}
\]

Suppose $\langle v^2 \rangle = \langle v_x^2 \rangle + \langle v_y^2 \rangle + \langle v_z^2 \rangle = 3 \langle v_x^2 \rangle = d \langle v_x^2 \rangle$.  
\[
\Longrightarrow P = \frac{mN}{dV} \langle v^2 \rangle = \frac{2}{d} \frac{1}{2} m \langle v^2 \rangle \frac{N}{V}
\]

\subsubsection*{An empirical gas law}
Now 
\[
P = \frac{N\tau}{V} \quad \, \text{(empirical)}
\]

\[
\Longrightarrow \frac{1}{2} m \langle v^2 \rangle = \frac{d}{2} \tau
\]

